\section{共同诉讼}
\begin{enumerate}
  \item 一方 \hz 双方为两人以上。
  \item 普通共同诉讼。\zd{诉讼标的同种类}。
        \begin{enumerate}
          \item \zd{若干个独立的可分之诉}。
          \item 合并审理（提高诉讼效率）必须经过当事人同意。
        \end{enumerate}
  \item 必要共同诉讼。\zd{诉讼标的同一}。\zd{一个事}。怎么区分，做题\zd{数诉讼标的的个数}。
        \begin{enumerate}
          \item 一个诉讼，必须合并审理。
          \item 关系紧密，因此\zd{一人行为，经其他共同诉讼人的同意}，对其他共同人发生法律效力。
          \item \zd{两个例外}。
                \begin{enumerate}
                  \item 赡养费的案件，所有的赡养义务作为共同被告。生活常理。
                  \item 遗产继承纠纷，部分继承人起诉的，应当通知其他继承作为共同原告参加诉。
                        不愿意参加诉讼又不明确表示放弃实体权利的，依然将其列为共同原告。
                \end{enumerate}
        \end{enumerate}
  \item 代表人诉讼。10 人以上， 代表 2 - 5 名。
        \begin{enumerate}
          \item 涉及到\zd{实体权利}的事项必须经过被代表的当事人\zd{全体同意}（民诉中只有一种同意）。
          \item 分类：起诉时人数众多一方当事人人数是否确定的。
                \begin{enumerate}
                  \item 人数确定的。代表人的产生。
                        \begin{enumerate}
                          \item 推选：全体的选全体（一致推举）；部分的选部分（存在派别）；
                          \item 选不出代表人。必要共同 \yjt 自己上；普通共同 \yjt 另行起诉。
                        \end{enumerate}
                  \item 人数不确定的。一定是普通共同诉讼。代表人的产生。
                        \begin{enumerate}
                          \item 推选同人数确定的；
                          \item 选不出代表人 \yjt 法院提出人选协商 \yjt 协商不成 \yjt 法院指定。
                          \item 就是不同意 \yjt 另行起诉（无意义）\yjt 法院下达裁定，适用之前的判决。
                                （\zd{预决效力}\footnote{避免类似案，不同判，激化社会矛盾。只能用于人数不确定的代表人诉讼。}）。
                        \end{enumerate}
                \end{enumerate}
        \end{enumerate}
\end{enumerate}
